% Separate, inactive extensions.
\section{Explicit Extensions: Not the Baseline}
\label{sec:p11-extensions}

\textbf{Separation rule.} Every object in this section is inactive in the
baseline. None enters its symbols, assumptions, equilibrium, propositions, or
empirical mapping. Activating any block would require a distinct model
specification with its own assumptions, proofs, and data justification.

\subsection{Extension 1: Smooth route choice}
\label{subsec:p11-smooth-choice}

If a future quantitative application requires smooth route shares, it may
replace the local deterministic utility comparison by
\[
  U_{ir}=W_{ir}+\epsilon_{ir},
  \qquad r\in\{I,E,T,A\}.
\]
Under an independently justified iid Type-I extreme-value distribution with
scale \(\sigma>0\), the extension-only probability would be
\begin{equation}
  P_{ir}^{\mathrm{logit}}
  =
  \frac{\exp(W_{ir}/\sigma)}
  {\sum_{\ell\in\{I,E,T,A\}}\exp(W_{i\ell}/\sigma)}.
  \label{eq:p11-logit-extension}
\end{equation}
This formula is not used by any baseline demand, cutoff, proposition, or
equilibrium equation. Activation requires route-share data, a quantitative
need, and shock-scale identification. Deterministic
\(r_i^*\) remains the theoretical baseline.

\subsection{Extension 2: Route-specific implementation}
\label{subsec:p11-route-realization}

If later evidence supports route-specific post-assignment implementation
probabilities, an extension may introduce \(\chi^I(q,m)\) and
\(\chi^E(q,m)\). These would replace neither the baseline route choice nor the
policy-invariant downstream probability without a new institutional model.
In particular,
\[
  M=1
  \ \not\Rightarrow\
  \chi^E(q,m)\uparrow
\]
by assumption. Activation requires an operational definition, route-aligned
data, and a causal justification for any policy dependence.

\subsection{Extension 3: Transfer-market microfoundation}
\label{subsec:p11-transfer}

A future adoption-cost, search, or bargaining block may microfound the outside
value \(T(q,m)\). No bargaining protocol, transfer price, adoption cost, or
counterparty choice is selected here. Until a separately specified design
specifies agents, information, surplus, outside options and timing,
\(T(q,m)\) remains the exogenous non-retained outside option of the baseline.

\subsection{Extension 4: Dynamic firm evolution}
\label{subsec:p11-dynamics}

A recursive model is admissible only when a future question and evidence
support a genuine state transition, such as accumulated manufacturing
capability, a persistent product portfolio, entry/exit, or transition from a
research-oriented developer to an integrated incumbent. No Bellman equation is
written here because no state vector or transition law has been specified. This
prevents static continuation-value accounting from being relabeled as a
dynamic foundation.

\subsection{Extension 5: Multi-CMO matching}
\label{subsec:p11-matching}

A richer search or matching model may be considered if match-level data
identify developers, suppliers, technology compatibility, search duration,
prices and match timing. No matching function, bargaining rule or congestion
externality is imposed here. The baseline continues to clear one aggregate
qualified-capacity market at \(p_m^*\).

\subsection{Extension 6: Research-versus-development allocation}
\label{subsec:p11-research-development}

An optional, separately specified block may use two controls: upstream research
\(x_i^R\) and project development \(x_i^D\). A resource or financing
constraint might link them. The block could be used to study
\(x_i^D\uparrow\) together with \(x_i^R\downarrow\). The block is not
implemented here: no constraint, objective, FOC, patent-production function,
or policy sign is specified.

Activation requires a separate decision because it adds a control dimension,
can shift the paper toward financial constraints, and overlaps with the
empirical mechanism studied by Shi Gu. The baseline retains one common
\(x_i\), commercialization organization, retained authorization,
manufacturing-capability matching, and CMO scarcity.

\subsection{Isolation conclusion}
\label{subsec:p11-isolation}

The six optional blocks are documented so that future work has explicit
insertion points. They remain inactive and cannot be cited as baseline
mechanisms or assumptions. Conversely, the baseline does not import this
file or depend on any extension-only symbol.
